% Generated draft: requires mathematical and source review.
\documentclass[11pt]{article}
\usepackage[margin=25mm]{geometry}
\usepackage{fontspec}
\setmainfont{Latin Modern Roman}
\usepackage{amsmath,amssymb,mathtools,unicode-math}
\setmathfont{Latin Modern Math}
\usepackage{enumitem,needspace,xurl,hyperref}
\setlength{\parindent}{0pt}
\setlength{\parskip}{5pt}
\setlength{\emergencystretch}{3em}
\newcommand{\N}{\mathbb{N}}
\newcommand{\Z}{\mathbb{Z}}
\newcommand{\Q}{\mathbb{Q}}
\newcommand{\R}{\mathbb{R}}
\newcommand{\C}{\mathbb{C}}
\newcommand{\F}{\mathbb{F}}
\newcommand{\A}{\mathbb{A}}
\newcommand{\PP}{\mathbb{P}}
\newcommand{\E}{\mathbb{E}}
\newcommand{\eps}{\varepsilon}
\newcommand{\dd}{\,\mathrm{d}}
\begin{document}
\section*{UnsolvedMath ID: 30004021}
\textbf{Source problem code:} catalogue entry
\subsection{Definitions and mathematical statement}
Write \(\mathfrak{sl}_n\) for the Lie algebra of trace-zero \(n\times n\) matrices over the ground field intended by the source; that field is not specified in the supplied record. A nilpotent element has Jordan partition \((3,2^2)=(3,2,2)\) if its Jordan blocks have sizes \(3,2,2\), and partition \((2^2)=(2,2)\) if its Jordan blocks have sizes \(2,2\). The notation \(W_k(\mathfrak g,f)\) denotes the source's W-algebra associated with the Lie algebra \(\mathfrak g\), level \(k\), and nilpotent element \(f\); its precise construction and conventions are not supplied.

Let \(f\in\mathfrak{sl}_7\) and \(f'\in\mathfrak{sl}_4\) be nilpotent elements with orbit partitions \((3,2^2)\) and \((2^2)\), respectively. The supplied conjecture asks whether
\[
W_{-7+7/3}(\mathfrak{sl}_7,f)
\cong
W_{-4+4/3}(\mathfrak{sl}_4,f').
\]
The levels are \(-14/3\) and \(-8/3\), respectively.

The conjecture is conditional on recovering the source's conventions: the supplied record does not specify the ground field, whether these are universal or simple W-algebras, or the category of isomorphism and the additional structures it must preserve. Consequently, the displayed identification retains the supplied target but does not resolve these missing specifications.
\subsection{Short English statement}
The conjecture asks whether the W-algebra associated with \(\mathfrak{sl}_7\), level \(-14/3\), and nilpotent Jordan blocks \((3,2,2)\) is isomorphic to the one associated with \(\mathfrak{sl}_4\), level \(-8/3\), and blocks \((2,2)\). The intended W-algebra construction and isomorphism conventions require source clarification.
\subsection{Sources}
\begin{itemize}
\item Supplied local UnsolvedMath registry record, permanent ID \(30004021\), title \emph{Exceptional Isomorphism of Two W-Algebras}, source problem code \texttt{catalogue entry}, in the Oberwolfach Reports Open Problems collection. No individual proposer is supplied; the proposed year is recorded as 2018.
\item Supplied report citation: \emph{Enveloping Algebras and Geometric Representation Theory} (2019), Oberwolfach Reports, pp.\ 3081--3146, DOI \texttt{10.4171/owr/2018/52}. Supplied source and literature-reference link, labeled ``Oberwolfach report 52'': \url{https://doi.org/10.4171/owr/2018/52}.
\item Literature assessment checked at \texttt{2026-08-21}: the supplied dated curation labels the problem open, reports that the report explicitly calls this identification a conjecture, and reports that no authoritative primary source located through its cutoff establishes or disproves this exact identification. This is supplied dated curation, not an independently verified literature assessment.
\item The registry supplies the statement-status label \texttt{corrected\_verified}, with an explanation that mathematical content was preserved while presentation or omitted definitions were repaired. This is a supplied editorial qualification, not independent verification of the mathematical statement or its conventions.
\end{itemize}
\end{document}
